%-------------------------------------------------------------------------------
%    SECTION TITLE
%-------------------------------------------------------------------------------
\cvsection{Research Experience}

%-------------------------------------------------------------------------------
%    CONTENT
%-------------------------------------------------------------------------------

\cvsubsection{Positions}

\begin{cventries}

    %----------------------------------------------------------
    \cventry
    {Vikram Sarabhai Space Centre, Indian Space Research Organisation} % Job title
    {Scientist/Engineer - SC} % Organization
    {Thiruvananthapuram, India} % Location
    {Oct 2024 -- Present} % Date(s)
    {}

\end{cventries}
% ----------------------------------------------------------

\cvsubsection{Projects}

\begin{cventries}

    \cventry
    {Advisor: Dr. Yuxiang Qin} % Job title
    {Synergy between future 21-cm experiments and physical cosmology: Effect of X-ray heating from Quasars on the high-redshift 21-cm signal} % Title
    {ANU-FRT Project, Canberra} % Location
    {June 2024 -- Aug. 2024} % Date(s)
    { % Description(s) of tasks/responsibilities
        As part of the Future Research Talent (FRT) program at the Australian National University, this project aimed at exploring the impact of AGNs on high redshift 21-cm signal by forward modeling the X-ray heating in semi-analytic galaxy formation and evolution model, MERAXES. The project involved training on the use of \texttt{MERAXES}, \texttt{21cmFAST} and \texttt{DRAGONS} to run cosmological simulations and analyze simulation outputs.
    }

    \vspace{1mm}

    \cventry
    {Advisor: Prof. R. Srianand} % Job title
    {Exploring growth and radiative properties of supermassive black holes at cosmic dawn using NINJA simulations} % Title
    {Masters Thesis, IUCAA Pune} % Location
    {Aug. 2023 -- May 2024} % Date(s)
    { % Description(s) of tasks/responsibilities
        This masters thesis project aimed to study the growth of supermassive black hole (SMBH) progenitors in cosmological simulations run using \texttt{MP-GADGET}. We investigated the evolution of emission spectrum due to accretion of gas onto SMBHs and estimate the contribution from growing seed black holes.
    }

    \vspace{1mm}

    \cventry
    {Advisor: Prof. R. Srianand} % Job title
    {Supermassive Black holes at Cosmic Dawn} % Title
    {VSP Project, IUCAA Pune} % Location
    {June 2023 -- Aug. 2023} % Date(s)
    { % Description(s) of tasks/responsibilities
        A 7-week funded summer project to understand cosmological simulations and the numerical methods used to simulate galaxy formation and astrophysical phenomena mainly focused on black holes and implementations of sub-grid physical models. I used generated datasets to study the evolutionary growth of black holes over time via accretion and mergers.
    }

    \vspace{1mm}

    \cventry
    {Advisor: Vivek Singh, Sudhir Sajan} % Job title
    {Weather Classification Using Images} % Title
    {Online Internship, Spartificial} % Location
    {Jan. 2023 -- March 2023} % Date(s)
    { % Description(s) of tasks/responsibilities
        In this online internship project, I present a potential weather classification method using Convolutional Neural Networks (CNN). It includes the construction of a CNN architecture constructed using \texttt{TensorFlow} and used VGG16 as a benchmark model. The performance metrics of the CNN were compared with that of the VGG16-based classifier on a publicly available weather dataset and showed that the CNN achieved comparable accuracy with the VGG16-based classifier.
    }

    \vspace{1mm}

    \cventry
    {Advisor: Prof. R. Srianand} % Job title
    {Studying the spectral emission from accreting supermassive black holes} % Title
    {Summer Internship, IUCAA Pune} % Location
    {June 2022 -- Aug 2022} % Date(s)
    { % Description(s) of tasks/responsibilities
        As part of a credited summer internship and training project, I reviewed the standard model of accretion onto compact objects via a thin $\alpha$-disk and simulated the spectrum from an accreting black hole and explored the characteristic timescales for black hole growth.
    }

    \vspace{1mm}

    \cventry
    {As part of academic coursework}
    {Mini-Projects}
    {IIST, Thiruvananthapuram}
    {Aug. 2022 -- May 2024}
    {
        \begin{cvitems}
        \item Estimating $\Lambda$CDM cosmological model parameters from Type Ia Supernovae data using Bayesian analysis
        \item Simulating the evolution and dynamics of the benchmark $\Lambda$-CDM cosmological model
        \item Clustering in galaxies from GAMA survey data using two-point correlation function
        \item Estimating star formation rates in galaxies using spectral indicators
        \item Studying the distribution of massive stars to understand Galactic morphology
        \item Analyzing HII regions around stars of different spectral classes
        \end{cvitems}
    }

    %---------------------------------------------------------
\end{cventries}
